\chapter{INTRODUCTION}

\ChapterSection{Introduction}

A rice farmer in the Terai plains photographs a cluster of brown lesions spreading across her
crop's leaves. She needs to know whether the damage is bacterial leaf blight, brown spot, or
blast---and she needs to know within hours, not days, because the wrong treatment applied too
late can mean the difference between a manageable loss and a ruined harvest. In Nepal, where
agriculture accounts for approximately one-quarter of national GDP and directly sustains the
majority of rural households \parencite{moaldAgricultureStats2080,nsoAgricultureCensus2021},
this scenario repeats across millions of smallholdings every growing season.

The gap between the farmer's need for timely diagnosis and the availability of expert
pathological assessment defines the core problem. Traditional diagnosis depends on extension
workers, agro-vet consultants, or laboratory confirmation---channels that are valuable but
geographically sparse, seasonally overloaded, and often unreachable at the moment a farmer
first observes symptoms in the field. The FAO estimates that plant pests and diseases can
destroy up to 40 percent of global crop production annually, with actual loss depending on crop,
location, season, and management practice \parencite{faoPlantHealthLosses}. Nepal's
topographic diversity, settlement patterns, and variable connectivity amplify this challenge
for smallholder farmers who may possess a smartphone and a crop image but lack immediate access
to a trained plant-health expert \parencite{khatriNepalDigitalAgri2023}.

Computer vision offers a partial technological response. Object detection can identify relevant
plant structures within complex field images, while vision-language models can generate
structured diagnostic interpretations from visual input. Mobile technology provides the
interface: the smartphone is the most realistic access point for field-level users in Nepal's
agricultural context \parencite{faoNepalDigitalAgriculture}. Yet connecting these capabilities
into a coherent, farmer-usable workflow requires engineering far beyond model training. A
practical diagnosis-support system must handle image capture, model inference, structured output
parsing, backend validation, advisory generation, local record keeping, offline resilience, and
uncertainty communication---each presenting its own design and implementation challenges.

Krishi Vaidya is an integrated final-year B.Sc.\ CSIT project that addresses this end-to-end
workflow. It composes four major subsystems: a YOLO-based plant-part detector, a Qwen-based
vision-language model for disease interpretation, an Express/Mongoose backend API for
orchestration and persistence, and a React Native/Expo mobile application for farmer-facing
interaction. The project is positioned within Nepal's digital agriculture direction as a
student-built prototype for crop-health support---not as a certified diagnostic authority, but
as an investigation of how trained model components can be integrated with application software
to serve a practical agricultural need.

\ChapterSection{Problem Statement}

Farmers require timely, understandable, and accessible support for crop disease diagnosis, but
expert assessment is seldom available at the field level when symptoms are first observed.
Existing digital approaches typically address isolated aspects of this problem: image
classification without advisory output, model training without deployment evaluation, or mobile
interfaces without offline resilience. A gap persists between what individual model components
can achieve in isolation and what a farmer-facing advisory system requires as an integrated
workflow.

The research problem addressed by this project is:

\begin{quote}
Can an integrated system combining object detection, vision-language interpretation, backend
orchestration, and an offline-aware mobile application be designed, implemented, and evaluated
to support image-based crop-health diagnosis and structured advisory output for the Nepalese
agricultural context?
\end{quote}

This question decomposes into three constituent challenges:

\ChapterSubsection{Visual detection challenge}

Crop images contain multiple plant structures, and disease-related reasoning often depends on
which part of the plant is visible. A system that ignores plant-part localization treats
complex field images as undifferentiated classification inputs. Krishi Vaidya addresses this
by training a YOLOv8 nano detector on a curated multi-crop plant-part dataset with eleven
agricultural object classes, then evaluating detection performance across thresholds and
quantized deployment variants.

\ChapterSubsection{Disease interpretation challenge}

Disease diagnosis requires more than detecting plant objects. The system must interpret visual
evidence and produce structured information---disease name, confidence, severity, treatment
guidance---that downstream application components can parse, validate, and display. Krishi
Vaidya addresses this through a LoRA-adapted VLM pipeline trained on structured image-text
records, with evaluation focused not only on classification accuracy but also on whether
generated outputs can be reliably parsed into canonical labels.

\ChapterSubsection{Mobile access and advisory workflow challenge}

Even accurate models are incomplete if farmers cannot access them through a practical interface.
The mobile application must support image capture, diagnosis submission, offline record keeping,
synchronization, crop management, and understandable result presentation. The backend must
protect the API, validate requests, orchestrate inference, store records, and return consistent
structured responses. Krishi Vaidya addresses this by implementing both the mobile and backend
layers as full application components rather than limiting the project to model training alone.

\ChapterSection{Objectives}

\ChapterSubsection{General objective}

To design, implement, and evaluate an integrated crop-health diagnosis-support system that uses
object detection, vision-language modelling, backend services, and a mobile application to
assist farmers with image-based disease interpretation and structured advisory output.

\ChapterSubsection{Specific objectives}

\begin{enumerate}[label=\arabic*.]
\item To investigate the viability of lightweight object detection for agricultural plant-part
localization by curating a multi-source dataset, training a YOLOv8 nano detector, and
evaluating detection performance across confidence thresholds and quantized deployment variants.
\item To evaluate LoRA-based fine-tuning of a vision-language model for structured disease
interpretation by preparing a Qwen-format dataset, training with quantized adapters, and
assessing both classification accuracy and output parseability.
\item To implement a layered backend API that orchestrates authentication, crop records, scan
handling, diagnosis inference, advisory generation, and validation within a documented
architectural pattern.
\item To implement an offline-aware mobile application that supports farmer-facing scan workflow,
bilingual output, local persistence, synchronization infrastructure, and crop-management
record keeping.
\item To integrate the subsystems into a documented diagnosis workflow and evaluate the
integrated path through runtime-verified contracts, generated architecture diagrams, and
completed distributed deployment evidence.
\end{enumerate}

\ChapterSection{Scope and Limitation}

The scope of this project encompasses the design, implementation, and manuscript documentation
of the Krishi Vaidya prototype across its four subsystems.

The \textbf{model scope} includes YOLO object detection for plant-part localization across
eleven agricultural classes and VLM-based disease interpretation for a subset of crop-disease
categories derived from the PlantVillage corpus. The \textbf{backend scope} includes the
implemented Express/Mongoose service architecture, scan and diagnosis workflow, advisory
response handling, authentication, security middleware, validation, and available test coverage.
The \textbf{mobile scope} includes the React Native/Expo application with scan workflow, local
SQLite persistence, offline-aware behavior, synchronization infrastructure, crop and task
records, bilingual interface, and on-device TFLite model assets.

The following boundary conditions apply:

\begin{enumerate}[label=\arabic*.]
\item The project does not claim to produce a certified, agronomically validated, or
government-approved diagnosis system. It is a research and engineering prototype suitable
for demonstrating integrated software and model development.
\item The YOLO dataset, though large and curated, exhibits significant class imbalance; some
classes have orders of magnitude fewer annotations than others.
\item The VLM pipeline is implemented and measurable, but current prediction metrics reveal
substantial parseability limitations that must be acknowledged as a central finding.
\item Due to time and budget constraints, automated backend test coverage prioritizes the core scan domain, while external API integrations and standard CRUD endpoints rely on extensive manual validation.
\item The mobile application's integration architecture, including its offline-to-online recovery mechanisms, has been thoroughly validated through physical-device testing and distributed runtime measurements.
\item The advisory output has not been validated by certified agricultural experts and must
not substitute for expert consultation.
\end{enumerate}

These limitations define the current state of the prototype and the specific directions for
future work, both of which are discussed in Chapter 6.

\ChapterSection{Development Methodology}

Krishi Vaidya follows an iterative software-development methodology organized around analysis,
design, implementation, and testing. The analysis phase identified the farmer-facing diagnosis
workflow, the model requirements for plant-part detection and disease interpretation, the backend
requirements for validation and orchestration, and the mobile requirements for offline-aware use.
The design phase translated those requirements into four cooperating subsystems: a YOLO detector,
a Qwen-based VLM service, an Express/Mongoose backend, and a React Native/Expo mobile application.

Implementation proceeded component by component. The YOLO pipeline was developed around dataset
curation, training, validation, threshold analysis, and quantized export. The VLM pipeline was
developed around PlantVillage-derived disease data, Qwen-format supervised fine-tuning export,
LoRA training, inference serving, and parseability evaluation. The backend was implemented as a
layered API with controllers, services, repositories, validation middleware, authentication, image
upload support, diagnosis orchestration, and advisory generation. The mobile application was
implemented around onboarding, crop records, scan capture, image editing, diagnosis display,
local SQLite persistence, cached advisory output, and synchronization infrastructure.

Testing and validation were performed at the component and system-boundary levels using the
artifacts available in the project workspace: YOLO validation outputs, VLM metric scripts and API
tests, backend Jest unit and integration tests, mobile inference tests, source-backed screenshots,
and documented integration contracts. The methodology is therefore evidence-first: claims in the
report are tied to project code, generated outputs, screenshots, or measured model results.

\ChapterSection{Report Organization}

Chapter 2 reviews the background study and literature related to digital agriculture, plant
disease detection, object detection, vision-language models, offline-first mobile architecture,
and ML system design. Chapter 3 presents system analysis, including requirements, feasibility,
and object-oriented analysis. Chapter 4 presents system design and algorithm details. Chapter 5
presents implementation, testing, and result analysis. Chapter 6 concludes the report and
provides future recommendations.
